\section{Watching It Changes It}
\label{sec:ch17-watching-it-changes-it}
\index{Advanced Request Tracing!and the plan cache|(}%

Chapter~\ref{ch:how-a-federated-query-actually-runs} named Advanced Request
Tracing once, in a sentence about why the sample router runs in development
mode, and never printed any of it. It is the router's answer to the request
timeline this book has been reading off HotChocolate consoles since
chapter~\ref{ch:the-life-of-a-request}, and it reports the phases you would
expect:

\begin{minted}[fontsize=\footnotesize,breaklines]{text}
parse_stats      duration_nanoseconds
normalize_stats  duration_nanoseconds
validate_stats   duration_nanoseconds
planner_stats    duration_nanoseconds
\end{minted}

Then a fetch tree, with every subgraph request and response body in it, which is
where section~\ref{sec:ch17-one-request-per-subgraph}'s representation counts
came from.

\subsection*{The line that spoils it}

\index{plan cache!skipped when you look at it}%
In the planner, before anything else happens:

\begin{minted}[fontsize=\footnotesize,breaklines]{go}
	// if we have tracing enabled or want to include a query plan in the response we always prepare a new plan
	// this is because in case of tracing, we're writing trace data to the plan
	// in case of including the query plan, we don't want to cache this additional overhead

	skipCache := options.TraceOptions.Enable || options.ExecutionOptions.IncludeQueryPlanInResponse
\end{minted}

The branch that follows returns before the cache is read, before it is written,
and before the single-flight group. The two request headers that let you see
what the planner did, \code{X-WG-Trace} and \code{X-WG-Include-Query-Plan}, are
the two that guarantee it did it.

\index{plan cache!the traced control}%
The comment gives a reason, and the reason is sound: trace data is written onto
fields of the plan object itself, so a traced request sharing a cached plan with
concurrent requests would be scribbling on something other requests are reading.
Given that design, bypassing the cache is the only safe thing to do. But sound
reasons still have consequences, and this one has a sharp consequence.

I ran it rather than trusting the read. One document, eight requests:

\begin{minted}[fontsize=\footnotesize,breaklines]{text}
untraced-1   planHit=false
untraced-2   planHit=true
untraced-3   planHit=true
traced-1     planHit=false
traced-2     planHit=false
planonly     planHit=false
untraced-4   planHit=true
untraced-5   planHit=true
\end{minted}

Every request is the same document with the same plan key. The first is a cold
miss. Then two hits. Then two traced requests, both of which miss, on a document
this router has planned twice already. Then the untraced requests afterwards hit
again, so a traced request does not poison the cache either; it simply is not
part of it.

\index{Advanced Request Tracing!reports only the cold number}%
So the planning duration in a trace is always the cold number. Two consecutive
traced requests for that document reported 832,666 and 1,083,310 nanoseconds of
planning, and both are measurements of work that would not have happened at all
without the header asking for them. Somebody tuning a router, reading
\code{planner\_stats}, and concluding that planning costs about a millisecond per
request would be wrong by however good their cache hit rate is, and the tool
that told them cannot tell them that.

The warm cost is not observable through the mechanism that reports the cold one,
by construction. What is observable is whether a request hit, and there are two
ways to see it. \code{planCacheHit} is an access-log expression like any other,
which is what this chapter's cases are built on and what I would put in a log on
a real deployment. The other is a debug setting that puts five cache results
into response headers, and it is the one to reach for at a keyboard:

\begin{minted}[fontsize=\footnotesize,breaklines]{text}
cold      plan=MISS normalization=MISS
warm-1    plan=HIT normalization=HIT
warm-2    plan=HIT normalization=HIT
traced    plan=MISS normalization=HIT
planonly  plan=MISS normalization=HIT
warm-3    plan=HIT normalization=HIT
\end{minted}

The two columns are \code{X-WG-Execution-Plan-Cache} and
\code{X-WG-Normalization-Cache}, shortened by the case that prints them.

\index{plan cache!the skip is the plan cache alone}%
That is the same finding arriving through a completely different mechanism,
which is the only reason I am confident in it, and it says one thing the access
log could not. On the traced request the normalisation cache still reports a
hit while the plan cache reports a miss. The header is not telling you that
tracing bypasses the pipeline; it is telling you that tracing bypasses exactly
one cache out of the several a request passes through, which is what the source
says and what a single \code{planHit=false} could not have distinguished from a
request being treated as new all the way down.

\subsection*{Seven caches and a default of false}

\index{router!nine caches, seven with metrics}%
The server struct holds nine caches and the router registers metrics for seven
of them by name: the plan cache, three around normalisation, one for persisted
operations, one for validation, and one for the operation hash. The two with no
metrics are the complexity cache and the slow-plan fallback. It publishes none
of the seven unless \code{graphql\_cache} is turned on under the Prometheus
block, and that setting defaults to false.

I am not sure that default is wrong, exactly. Cache metrics are cardinality and
somebody decided you should opt in. What I am sure of is the effect it had on
this book: four chapters asked a question about a cache, a fifth had the
endpoint in front of it, and the answer arrived only when somebody stopped
looking for a counter and asked the router to log its own key. A router whose
caches are invisible is a router you cannot size, and sizing it is a thing the
configuration expects you to do, since the plan cache takes an entry count.

\code{router/config.yaml} in the companion repository carries the block as a
comment with this measurement beside it, in the same spirit as the WebSocket
block chapter~\ref{ch:identity-and-authorization} left there: a setting Mosaic
does not ship, recorded where somebody operating this graph would look for it.
Chapter~\ref{ch:observability} owns observability and is where it stops being a
comment.

\medskip
\index{router!what reading it was worth}%
Reading two processes nobody on your team wrote is a strange way to spend a
chapter, and it has paid differently from
chapter~\ref{ch:inside-hotchocolate}'s. That chapter mostly confirmed an
executor and found its surprises attached to lines somebody had written for
another reason. This one found that the composer blames a route rather than a
fault, that a plan is shared by more documents than I would have guessed and
split by fewer things than I would have accepted, and that the instrument
reports a number the running system never pays. None of the three is a bug and
all three change what you would do.

The next chapter goes back out to the wire, on the hop this one has not looked
at: chapter~\ref{ch:the-wire-completely} owns everything between a client and
the router, where this chapter has been standing entirely on the other side of
it.

\index{Advanced Request Tracing!and the plan cache|)}%
